\begin{claim}{D.2.26}{Proved}{The exact quadratic test for a simultaneous first cone}
All groups below have coefficients in $\FF_2$.
For every three-connected CW space $C$, Pontryagin multiplication
induces a natural injection
\[
 \frac{(H_4C\otimes H_5C)\oplus(H_5C\otimes H_4C)}
 {\im(\overline\Delta:H_9C\longrightarrow
                (H_4C\otimes H_5C)\oplus(H_5C\otimes H_4C))}
 \lhook\joinrel\longrightarrow H_7(\Omega C).
 \tag{D.2.26a}
\]
Here the factors are identified with $H_3(\Omega C)$ and
$H_4(\Omega C)$ by the canonical homology desuspensions.

Put $Y=\Sigma^2K(\FF_2,2)$, and let $u\in H^4Y$ be the
bottom class and $v=\operatorname{Sq}^1u\in H^5Y$.
Write $a,b$ for the degree-three and degree-four letters in
$H_*(\Omega Y)$, and $\kappa=ab+ba$.
Let $h=\bigvee_i h_i:S=\bigvee_i S_i\to Y$, where every $S_i$
is two-connected and $h_i$ is zero on $H_4,H_5$, and set
$C=Y\cup_h CS$. The old commutator survives in $H_7(\Omega C)$
exactly when its tensor $a_4\otimes b_5+b_5\otimes a_4$ is
outside the image of the diagonal in \textup{(D.2.26a)}.
If $\alpha\in H^4C$ lifts $u$, the condition
\[
 \alpha\operatorname{Sq}^1\alpha=0
 \tag{D.2.26b}
\]
is independent of the lift and implies this survival. It is
equivalent to survival when $H_4C,H_5C$ are precisely the old
one-dimensional groups.

There is an exact simultaneous test for \textup{(D.2.26b)}.
Fix singular cocycles $a^4,b^5$ representing $u,v$ and a cochain
$c^8$ on $Y$ with $dc^8=a^4\smile b^5$. For each $i$ choose
$p_i^3$ with $dp_i^3=h_i^*a^4$, and put
\[
 \theta_i=[h_i^*c^8-p_i^3\smile h_i^*b^5]\in H^8S_i.
 \tag{D.2.26c}
\]
Then \textup{(D.2.26b)} holds if and only if there is one
$\ell\in H^8Y$ such that $\theta_i=h_i^*\ell$ for every $i$.
Distinct choices of $\ell$ for the separate cones do not suffice.

If, moreover, $S_i=\Sigma R_i$ with $R_i$ two-connected and
$h_i$ also zero on $H_6$, define $L_i:H_8S_i\to\FF_2$ by the
ordinary adjunct $\bar h_i:R_i\to\Omega Y$ and the fixed
James decomposition:
\[
 (\bar h_i)_*z=(h_i)_*(\sigma z)+L_i(\sigma z)\kappa,
 \quad
 H_7(\Omega Y)=H_8Y\oplus\FF_2\{ab,ba\}.
 \tag{D.2.26d}
\]
The first term lies in the word-length-one summand. In this
suspended-source case the exact necessary and sufficient condition
for survival is
\[
 L=\bigoplus_iL_i\text{ vanishes on }
 \ker\!\left(\bigoplus_iH_8S_i\xrightarrow{h_*}H_8Y\right),
 \quad\text{equivalently }L_i=\ell\circ(h_i)_*\text{ for one }\ell.
 \tag{D.2.26e}
\]
The assertions apply to arbitrary small wedges and CW sources.
They are criteria for the first cone, and assert neither that
all the attaching maps satisfy them nor persistence at later stages.
\end{claim}
\begin{proof}
For a finite-type three-connected $C$, take a minimal
Adams--Hilton model $(T(V),d)$ over $\FF_2$.
Here $V=s^{-1}\widetilde H_*C$ starts in degree three, $d$ has
no linear part, and its quadratic part is the desuspended reduced
diagonal; these are the exact assertions of
\cite[\S3.1, p.300]{FGHHalperinLemaire1995}.
The degree-seven length-two words are $V_3V_4+V_4V_3$ and are
cycles. Words of length at least three start in degree nine.
A degree-eight length-two word has letters of degree at most five,
whose differentials vanish. Hence all boundaries in this
degree-seven subspace are exactly $d(V_8)$, which is quadratic
and equals the diagonal of $H_9C$. This proves
\textup{(D.2.26a)}. Its definition by the actual Pontryagin product
proves naturality and rules out an unresolved filtration extension.

For general $C$, use a CW model with cells only in dimensions zero
and at least four and take its filtered finite subcomplexes.
All are three-connected. Homology, field tensor products, the
diagonal and its cokernel commute with this union. So does loop
homology: a finite singular cycle or bounding chain in the loop
space has compact adjoint and lies in a finite subcomplex.
The natural injections therefore pass to the filtered colimit.

The cone $C$ is three-connected, and the old $H_4,H_5$ inject
by the homology cofiber sequence. Naturality gives the exact tensor
criterion. If \textup{(D.2.26b)} holds, the functional
$\alpha\otimes\operatorname{Sq}^1\alpha$ annihilates the
diagonal image and evaluates to one on the old tensor, proving
survival. If the two groups are one-dimensional, cocommutativity
makes that tensor span every possible diagonal image, proving
the converse under the stated extra hypothesis.

The ordinary coaction $C\to C\vee\Sigma S$ proves that
connecting-image classes annihilate positive cohomology of $C$:
projection to $C$ is homotopic to the identity, and pullback from
$\Sigma S$ is the connecting homomorphism, while products of
positive classes on the separate wedge summands are zero.
Any two lifts differ by a connecting-image class. Thus products
of lifts of fixed positive classes are independent of the lifts.

Let $\delta:H^8S\to H^9C$ be the connecting homomorphism.
We prove using genuine singular cup products that, for lifts
$\alpha,\beta$ of $u,v$,
\[
 \alpha\beta=\delta((\theta_i)_i).
 \tag{D.2.26f}
\]
First \textup{(D.2.26c)} is a cocycle by Leibniz. Changing $p_i$
by a closed cochain changes it by a product with $h_i^*b^5$,
which is exact, so its class is unchanged.

Use a mapping-cylinder cone cover $U,V$, with $U\simeq Y$,
$V$ contractible and $U\cap V$ a collar of $S$.
The small singular chains $C_*(U)+C_*(V)$ include
quasi-isomorphically in $C_*(C)$ by subdivision and retain the
actual Alexander--Whitney cup: faces of each small simplex stay
in the same cover member. Their dual has the exact sequence
\[
 0\to C^*_{\rm small}(C)\to C^*(U)\oplus C^*(V)
       \to C^*(U\cap V)\to0.
 \tag{D.2.26g}
\]
The last map is the difference of restrictions and is onto by
extension of cochains on singular simplices.
Choose global cocycles for $\alpha,\beta$ restricting exactly
to the pullbacks of $a^4,b^5$ on $U$; extend correcting cochains
from $U$ to achieve these representatives. On $V$ write
$\alpha=dP^3$. Their cup product has primitive $c^8$ on $U$
and primitive $P^3\smile\beta$ on $V$. The difference on the
collar, restricted to $S$, is precisely
$h^*c^8-p^3\smile h^*b^5$. The connecting construction in
\textup{(D.2.26g)} gives \textup{(D.2.26f)}.
No multiplication on an abstract mapping-cone DGA is assumed.

Take $\beta=\operatorname{Sq}^1\alpha$. Since
$H^8S=\prod_iH^8S_i$, the cofiber exact sequence now gives
\[
 \alpha\operatorname{Sq}^1\alpha=0
 \quad\Longleftrightarrow\quad
 (\theta_i)_i\in\im(H^8Y\xrightarrow{h^*}\prod_iH^8S_i).
 \tag{D.2.26h}
\]
This is exactly the common-functional assertion, by universal
coefficients over a field. It includes cancellations between
different summands of $H_8S$. For example, two input vectors
with the same nonzero target image but different functional
values are harmless to the separate kernel tests and violate
the simultaneous one. This is a linear-algebra observation,
not a claimed realization by attaching maps in this problem.

For the last assertion put $V=\widetilde H_*(\Sigma B)$,
$W=\widetilde H_*(\bigvee_iR_i)$. The James identification
gives $H_*(\Omega Y)=T(V)$ with primitive letters, since
$\Sigma B$ is itself a suspension. In degrees three, four and
five, homology suspension is an isomorphism; hence the low
vanishing hypotheses give $(\bar h)_*W_3=(\bar h)_*W_4
=(\bar h)_*W_5=0$. Coalgebra naturality in degree seven
then forces the $ab$ and $ba$ coefficients to agree, proving
that \textup{(D.2.26d)} is well-defined.

Both source and target are AH-suspensions. The opened relative
cone construction in the proof of
\cite[Proposition 7.4, p.309]{FGHHalperinLemaire1995} gives
\[
 AH(C)=(T(V\oplus sW),D),\qquad
 D|_V=0,\qquad D(sw)=(\bar h)_*w.
 \tag{D.2.26i}
\]
Old generators start in degree three, new ones in degree four.
The differentials of new generators in degrees four, five and
six vanish by the preceding low vanishing. Every decomposable
degree-eight chain therefore has zero differential: its factors
have these degrees or are old generators; a degree-seven new
generator would require a degree-one companion, and words of
length three start in degree nine. Thus all degree-seven
boundaries are exactly $D(sW_7)$. The old $\kappa$ is a boundary
if and only if a finite sum of the vectors in
\textup{(D.2.26d)} has zero linear component and nonzero
$\kappa$ component. This is precisely the negation of
\textup{(D.2.26e)}. In particular new $H_5C$ classes permitted
by a once-suspended source create no omitted boundary here.

In fact the two normalizations agree modulo one common pullback:
\[
 (\theta_i)_i-L\in\im(h^*:H^8Y\to H^8S).
 \tag{D.2.26j}
\]
For $z\in\ker h_*$ choose a lift $\widetilde z\in H_9C$.
After cancelling the linear differential in
\textup{(D.2.26i)}, its quadratic differential pairs with
$\alpha\otimes\beta$ as $L(z)$, by the reduced-diagonal
identification of \cite[\S3.1]{FGHHalperinLemaire1995}.
But \textup{(D.2.26f)} gives the same pairing as
$\langle\delta\theta,\widetilde z\rangle=\theta(z)$.
Thus $\theta-L$ vanishes on the whole kernel and factors
through $h_*$, proving \textup{(D.2.26j)}.

The cited paper permits characteristic two, with a finite-type
convention. For an arbitrary wedge, every input vector or
boundary witness uses finitely many source cells and summands.
Choose finite subpairs of the mapping-cylinder cone, enlarging
the finite target to contain chains witnessing the finitely many
required low homology vanishings. The finite-type relative
construction proves the same criterion for each finite witness,
and passage to their filtered union proves it as stated.
Finally a functional on the source sum vanishes on $\ker h_*$
exactly when it factors through $\im h_*$; extend it linearly to
$H_8Y$ to obtain the single $\ell$. No finite-dimensional duality
or countability assumption enters this step.
\end{proof}
